\documentclass[a4paper,11pt]{ltjsarticle}
\newcommand{\bookletname}{模範解答}
% LuaLaTeX / UTF-8  第2回模試 共通スタイル（統計数理）
\usepackage{amsmath,amssymb,bm}
\usepackage[margin=23mm,top=22mm,bottom=22mm]{geometry}
\usepackage{enumitem,booktabs,array,multirow}
\usepackage{fancyhdr}
\usepackage[hidelinks,unicode]{hyperref}
\setlength{\parindent}{1\zw}
\setlength{\parskip}{0.35em}
\setlength{\headheight}{17pt}
\pagestyle{fancy}
\fancyhf{}
\fancyhead[L]{\small 第2回模試\quad 統計数理}
\fancyhead[R]{\small\bookletname}
\fancyfoot[C]{\small\thepage}
\renewcommand{\headrulewidth}{0.35pt}
\renewcommand{\footrulewidth}{0pt}
\setlist[enumerate,1]{label=〔\arabic*〕,leftmargin=3.5em,labelsep=0.5em,itemsep=0.9em,topsep=0.7em}
\newcommand{\E}{\operatorname{E}}
\newcommand{\Var}{\operatorname{Var}}
\newcommand{\Cov}{\operatorname{Cov}}
\newcommand{\Corr}{\operatorname{Corr}}
\newcommand{\Prob}{\operatorname{P}}
\newcommand{\ind}{\mathbf{1}}
\newcommand{\qhead}[1]{\section*{問#1}}
\newcommand{\ans}[1]{\subsection*{〔#1〕}}
\newcommand{\points}[1]{\hfill\textnormal{（#1点）}}
\newcommand{\continued}[1]{\noindent\textbf{問#1（続き）}\par\medskip}
\hypersetup{pdfauthor={Claude Fable 5.2},pdfsubject={統計検定1級対策・自主学習用模擬試験}}

\setlength{\parskip}{0.2em}
\hypersetup{pdftitle={第2回模試 統計数理 模範解答}}
\begin{document}
\setlength{\abovedisplayskip}{6pt plus 2pt minus 2pt}
\setlength{\belowdisplayskip}{6pt plus 2pt minus 2pt}
\noindent{\small 付表番号・数値は「統計検定1級\_付表数値表」に対応する。数値計算は四則演算・平方根・付表だけで再現できる形で示し，付表から引いた値を明記している。}

\qhead{1}
\ans{1}
$u\mapsto u/(1-u)$ は $(0,1)$ で狭義単調増加で，値域は $(0,\infty)$ である。したがって $X$ は実数全体に値をとり，任意の実数 $x$ に対して
\[
F(x)=\Prob\!\left(\frac{U}{1-U}\le e^x\right)=\Prob\!\left(U\le\frac{e^x}{1+e^x}\right)
=\frac{e^x}{1+e^x}=\boxed{\frac{1}{1+e^{-x}}}.
\]
最後から2番目の等号は $U$ が $(0,1)$ 上の一様分布に従い，$0<e^x/(1+e^x)<1$ であることによる。微分して
\[
f(x)=F'(x)=\boxed{\frac{e^{-x}}{(1+e^{-x})^2}}\qquad(-\infty<x<\infty).
\]
$1-F(x)=e^{-x}/(1+e^{-x})$ であるから，
\[
F(x)\{1-F(x)\}=\frac{1}{1+e^{-x}}\cdot\frac{e^{-x}}{1+e^{-x}}=f(x).
\]
また，$f(-x)$ の分母・分子に $e^{-2x}$ を掛けると
\[
f(-x)=\frac{e^{x}}{(1+e^{x})^2}=\frac{e^{-x}}{(e^{-x}+1)^2}=f(x).
\]

\ans{2}
上側25\%点 $q$ は $F(q)=3/4$ を満たす。$1+e^{-q}=4/3$ より $e^{-q}=1/3$，すなわち $q=\log3$ である。
付表5より $\log_{10}3=0.4771$ であるから
\[
q=\log3=2.3026\times0.4771=1.0986\simeq\boxed{1.099}.
\]
次に，$f$ は偶関数であるから $\E[|X|]=2\int_0^\infty xf(x)\,dx$ である。
$x\ge0$ で $0\le1-F(x)=e^{-x}/(1+e^{-x})\le e^{-x}$ なので $x\{1-F(x)\}\to0$（$x\to\infty$）となり，部分積分により
\begin{align*}
\int_0^\infty xf(x)\,dx&=\Bigl[-x\{1-F(x)\}\Bigr]_0^\infty+\int_0^\infty\{1-F(x)\}\,dx\\
&=\int_0^\infty\frac{e^{-x}}{1+e^{-x}}\,dx
=\Bigl[-\log(1+e^{-x})\Bigr]_0^\infty=\log2 .
\end{align*}
よって
\[
\E[|X|]=\boxed{2\log2}\simeq2\times2.3026\times0.3010=1.386
\]
（付表5の $\log_{10}2=0.3010$ を用いた）。なお，$\E[|X|]<\infty$ と対称性から $\E[X]=0$ である。

\ans{3}
$Y_1,Y_2>0$ なので $W$ は実数全体に値をとる。任意の実数 $w$ に対し，$Y_2=y$ で条件付けて $Y_1$ と $Y_2$ の独立性を用いると，
\[
\Prob(W\le w)=\Prob(Y_1\le e^wY_2)=\int_0^\infty\Prob(Y_1\le e^wy)\,e^{-y}\,dy
=\int_0^\infty\bigl(1-e^{-e^wy}\bigr)e^{-y}\,dy .
\]
$\int_0^\infty e^{-(1+e^w)y}\,dy=1/(1+e^w)$ であるから
\[
\Prob(W\le w)=1-\frac{1}{1+e^w}=\frac{e^w}{1+e^w}=\frac{1}{1+e^{-w}}=F(w).
\]
分布関数が一致するので，$W$ は $X$ と同じ分布（標準ロジスティック分布）に従う。

\ans{4}
$e^{tX}=\{U/(1-U)\}^t=U^t(1-U)^{-t}$ であり，$U$ の密度は $(0,1)$ 上で1であるから
\[
M(t)=\int_0^1u^{t}(1-u)^{-t}\,du=\int_0^1u^{(1+t)-1}(1-u)^{(1-t)-1}\,du .
\]
被積分関数は正で，この積分が有限となるのは $1+t>0$ かつ $1-t>0$，すなわち $|t|<1$ のときである。このとき与えられた公式で $a=1+t,\ b=1-t$ とおくと，$\Gamma(2)=1$ より
\[
M(t)=\frac{\Gamma(1+t)\Gamma(1-t)}{\Gamma(2)}=\Gamma(1+t)\,\Gamma(1-t)\qquad(|t|<1).
\]
次に分散を求める。与えられた等式より $0<|t|<1$ で $M(t)=\pi t/\sin(\pi t)$ である。
$y=\pi t$ とおき，$\sin y=y-y^3/6+O(y^5)$（$y\to0$）を用いると
\[
\frac{y}{\sin y}=\frac{1}{1-y^2/6+O(y^4)}=1+\frac{y^2}{6}+O(y^4),
\qquad\text{すなわち}\qquad
M(t)=1+\frac{\pi^2}{6}t^2+O(t^4)\quad(t\to0).
\]
（$t=0$ でも $M(0)=1$ であり，この展開式は原点の近傍全体で成り立つ。）
積率母関数が原点の近傍で存在するので，$M(t)=1+\E[X]t+\E[X^2]t^2/2+\cdots$ と展開できる。係数を比較して
$\E[X]=0$，$\E[X^2]/2=\pi^2/6$ を得る。したがって
\[
\Var(X)=\E[X^2]-(\E[X])^2=\boxed{\frac{\pi^2}{3}}\ (\simeq3.290).
\]

\ans{5}
(a) $X_1,\ldots,X_n$ は独立に分布関数 $F$ をもつので，任意の実数 $x$ に対して
\[
\Prob(M_n-\log n\le x)=\{F(x+\log n)\}^n=\left(1+e^{-x-\log n}\right)^{-n}
=\left(1+\frac{e^{-x}}{n}\right)^{-n}.
\]
$(1+c/n)^n\to e^c$（$n\to\infty$）であるから
\[
\lim_{n\to\infty}\Prob(M_n-\log n\le x)=\boxed{\exp(-e^{-x})}.
\]
（これは連続な分布関数であり，$M_n-\log n$ はこの分布関数をもつ分布（ガンベル分布）に分布収束する。）

(b) $M_2$ の分布関数は $F(x)^2$，密度は $2F(x)f(x)$ である。$u=F(x)$ と置換すると $du=f(x)\,dx$，$x=\log\{u/(1-u)\}$ であるから
\[
\E[M_2]=\int_{-\infty}^{\infty}x\cdot2F(x)f(x)\,dx
=2\int_0^1u\log u\,du-2\int_0^1u\log(1-u)\,du .
\]
部分積分により $\displaystyle\int_0^1u\log u\,du=\Bigl[\tfrac{u^2}{2}\log u\Bigr]_0^1-\int_0^1\frac u2\,du=-\frac14$。
また $v=1-u$ とおくと
\[
\int_0^1u\log(1-u)\,du=\int_0^1(1-v)\log v\,dv=\int_0^1\log v\,dv-\int_0^1v\log v\,dv=-1+\frac14=-\frac34 .
\]
よって
\[
\E[M_2]=2\left(-\frac14\right)-2\left(-\frac34\right)=\boxed{1}.
\]

(c) $M_n$ の分布関数は $F(x)^n$，密度は $nF(x)^{n-1}f(x)$ である。(b) と同じ置換 $u=F(x)$ により
\[
\E[M_n]=n\int_0^1u^{n-1}\log u\,du-n\int_0^1u^{n-1}\log(1-u)\,du .
\]
第1項は，部分積分により
\[
n\int_0^1u^{n-1}\log u\,du=\Bigl[u^n\log u\Bigr]_0^1-\int_0^1u^{n-1}\,du=-\frac1n .
\]
第2項は，$nu^{n-1}$ の原始関数として $u^n-1$ をとって部分積分する。$u\to1$ のとき $(u^n-1)\log(1-u)\to0$ であるから
\begin{align*}
n\int_0^1u^{n-1}\log(1-u)\,du&=\Bigl[(u^n-1)\log(1-u)\Bigr]_0^1-\int_0^1\frac{1-u^n}{1-u}\,du\\
&=-\int_0^1\bigl(1+u+\cdots+u^{n-1}\bigr)du=-\sum_{k=1}^{n}\frac1k .
\end{align*}
よって
\[
\E[M_n]=-\frac1n+\sum_{k=1}^{n}\frac1k=\sum_{k=1}^{n-1}\frac1k .
\]
（$n=2$ のとき右辺は1で，(b) と一致する。）
\newpage

\qhead{2}
\ans{1}
確率モデルは「$X_1,\ldots,X_n$ が独立に $N(\mu,c^2\mu^2)$ に従う（$\mu>0$ は未知，$c>0$ は既知）」である。1個の観測値の確率密度関数は
\[
f(x;\mu)=\frac{1}{\sqrt{2\pi}\,c\mu}\exp\!\left\{-\frac{(x-\mu)^2}{2c^2\mu^2}\right\}\qquad(-\infty<x<\infty)
\]
であるから，観測値 $x_1,\ldots,x_n$ に対する尤度関数は
\[
L(\mu)=\prod_{i=1}^nf(x_i;\mu)=(2\pi c^2\mu^2)^{-n/2}\exp\!\left\{-\frac{1}{2c^2\mu^2}\sum_{i=1}^n(x_i-\mu)^2\right\}\qquad(\mu>0).
\]
$\sum_i(x_i-\mu)^2=n(t-2\mu\bar x+\mu^2)$（$\bar x,t$ は $\bar X,T$ の実現値）であるから
\[
\boxed{L(\mu)=(2\pi c^2\mu^2)^{-n/2}\exp\!\left\{-\frac{n(t-2\mu\bar x+\mu^2)}{2c^2\mu^2}\right\}}
\]
となる。同時密度が $(\bar x,t)$ と $\mu$ だけの関数（$\times$ 定数1）に分解されるので，分解定理（因子分解定理）により $(\bar X,T)$ は $\mu$ の十分統計量である。

\ans{2}
〔1〕の尤度から，対数尤度は
\[
\ell(\mu)=\log L(\mu)=-n\log\mu-\frac{n}{2c^2}\left(\frac{t}{\mu^2}-\frac{2\bar x}{\mu}+1\right)+\text{定数}.
\]
$\mu$ で微分すると（スコア関数）
\[
\ell'(\mu)=-\frac n\mu+\frac{n}{c^2}\left(\frac{t}{\mu^3}-\frac{\bar x}{\mu^2}\right)
=-\frac{n}{c^2\mu^3}\bigl(c^2\mu^2+\bar x\mu-t\bigr).
\]
2次関数 $q(\mu)=c^2\mu^2+\bar x\mu-t$ は $\mu^2$ の係数が正で，確率1で $t>0$ より $q(0)=-t<0$ であるから，
正の零点をちょうど1つもつ。それを $\hat\mu$ とすると
\[
\hat\mu=\frac{-\bar x+\sqrt{\bar x^2+4c^2t}}{2c^2}.
\]
$0<\mu<\hat\mu$ では $q(\mu)<0$ より $\ell'(\mu)>0$，$\mu>\hat\mu$ では $q(\mu)>0$ より $\ell'(\mu)<0$ である。
よって $\ell(\mu)$ は $\mu=\hat\mu$ で最大となる。最尤推定量は
\[
\boxed{\hat\mu=\frac{-\bar X+\sqrt{\bar X^2+4c^2T}}{2c^2}}.
\]

\ans{3}
1個の観測値の対数密度は $\log f(x;\mu)=-\log\mu-\dfrac{(x-\mu)^2}{2c^2\mu^2}+\text{定数}$ であり，
\[
\frac{\partial}{\partial\mu}\log f(x;\mu)=-\frac1\mu+\frac{x-\mu}{c^2\mu^2}+\frac{(x-\mu)^2}{c^2\mu^3}.
\]
$Z=(X_1-\mu)/(c\mu)$ とおくと $Z\sim N(0,1)$ であり，スコアは
\[
\frac{\partial}{\partial\mu}\log f(X_1;\mu)=\frac1\mu\left(\frac Zc+Z^2-1\right)
\]
と書ける。$\E[Z]=0$，$\E[Z^2]=1$ より期待値は0である。$\Var(Z)=1$，$\Var(Z^2)=\E[Z^4]-1=2$，$\Cov(Z,Z^2)=\E[Z^3]=0$ であるから
\[
I(\mu)=\Var\!\left(\frac{\partial}{\partial\mu}\log f(X_1;\mu)\right)
=\frac{1}{\mu^2}\left(\frac{1}{c^2}+2\right)=\boxed{\frac{1+2c^2}{c^2\mu^2}}.
\]
次に，$\E[\bar X]=\mu$ より $\bar X$ は不偏であり，$\Var(\bar X)=c^2\mu^2/n$ である。
一方，クラメール・ラオの不等式による $\mu$ の不偏推定量の分散の下限は
\[
\frac{1}{nI(\mu)}=\frac{c^2\mu^2}{n(1+2c^2)}
\]
である。$c>0$ より $1+2c^2>1$ なので
\[
\Var(\bar X)=\frac{c^2\mu^2}{n}>\frac{c^2\mu^2}{n(1+2c^2)}
\]
となり，$\bar X$ の分散は下限より大きい（$\bar X$ の効率は $1/(1+2c^2)<1$）。

\ans{4}
$\E[\tilde\mu]=\E[S]/(a_nc)=a_nc\mu/(a_nc)=\mu$ であるから，$\tilde\mu$ は不偏である。
$S^2$ は分散 $c^2\mu^2$ の不偏推定量なので $\E[S^2]=c^2\mu^2$ であり，
\[
\Var(S)=\E[S^2]-(\E[S])^2=c^2\mu^2(1-a_n^2),
\qquad
\Var(\tilde\mu)=\frac{\Var(S)}{a_n^2c^2}=\boxed{\frac{(1-a_n^2)\mu^2}{a_n^2}}.
\]
正規分布からの無作為標本では $\bar X$ と $S^2$ は独立であるから，$\bar X$ と $\tilde\mu$ も独立である。
$V_1=\Var(\bar X)=c^2\mu^2/n$，$V_2=\Var(\tilde\mu)$ とおくと
\[
\Var(\hat\mu_w)=w^2V_1+(1-w)^2V_2 .
\]
これは $w$ の下に凸な2次関数で，$w$ で微分して0とおくと $2wV_1-2(1-w)V_2=0$，すなわち
\[
w^*=\frac{V_2}{V_1+V_2}=\boxed{\frac{n(1-a_n^2)}{c^2a_n^2+n(1-a_n^2)}}
\]
（$0<w^*<1$ を満たし，$\mu$ に依存しない）。このとき
\[
V^*=\frac{V_1V_2}{V_1+V_2}=\boxed{\frac{c^2(1-a_n^2)\mu^2}{c^2a_n^2+n(1-a_n^2)}}.
\]

\ans{5}
(a) $S^2=\dfrac{n}{n-1}(T-\bar X^2)$ であるから，$S$ は $(\bar X,T)$ の関数である。そこで
\[
\boxed{h(\bar X,T)=\bar X-\frac{S}{a_nc}=\bar X-\frac{1}{a_nc}\sqrt{\frac{n}{n-1}\,(T-\bar X^2)}}
\]
とおくと，〔4〕よりすべての $\mu>0$ で $\E[h(\bar X,T)]=\mu-\mu=0$ である。
一方，$\bar X$ と $S$ は独立な連続型確率変数であるから $\Prob(\bar X=S/(a_nc))=0$ であり，$\Prob(h=0)=0<1$ である。
（したがって十分統計量 $(\bar X,T)$ は完備ではない。別の例として，$\E[T]=(1+c^2)\mu^2$，$\E[\bar X^2]=(1+c^2/n)\mu^2$ より
$h=T/(1+c^2)-n\bar X^2/(n+c^2)$ も条件を満たす。）

(b) $1-a_n^2=\dfrac{1}{2n}+O(n^{-2})$ より $n(1-a_n^2)\to\dfrac12$，$a_n^2\to1$ であるから
\[
nV^*=\frac{c^2\mu^2\cdot n(1-a_n^2)}{c^2a_n^2+n(1-a_n^2)}
\longrightarrow\frac{c^2\mu^2\cdot\frac12}{c^2+\frac12}=\boxed{\frac{c^2\mu^2}{1+2c^2}}=\frac{1}{I(\mu)}.
\]
すなわち，$\bar X$ と $S$ に基づく2つの不偏推定量の最良の線形結合 $\hat\mu_{w^*}$ は，$\bar X$ 単独とは異なり，
$n\to\infty$ でクラメール・ラオの下限 $1/\{nI(\mu)\}$ を漸近的に達成する。
\newpage

\qhead{3}
\ans{1}
同順位がないので，$(R_1,\ldots,R_n)$ と $(Q_1,\ldots,Q_n)$ はいずれも $1,\ldots,n$ の並べ替えである。したがって
\[
\bar R=\bar Q=\frac1n\cdot\frac{n(n+1)}2=\frac{n+1}2,\qquad
\sum_{i=1}^nR_i^2=\sum_{i=1}^nQ_i^2=\frac{n(n+1)(2n+1)}6
\]
であり，
\begin{align*}
A:=\sum_{i=1}^n(R_i-\bar R)^2&=\frac{n(n+1)(2n+1)}6-\frac{n(n+1)^2}4\\
&=\frac{n(n+1)\{2(2n+1)-3(n+1)\}}{12}=\frac{n(n^2-1)}{12}.
\end{align*}
同様に $\sum_i(Q_i-\bar Q)^2=A$ である。$\bar R=\bar Q$ より $R_i-Q_i=(R_i-\bar R)-(Q_i-\bar Q)$ なので，
\[
D=\sum_i(R_i-\bar R)^2+\sum_i(Q_i-\bar Q)^2-2\sum_i(R_i-\bar R)(Q_i-\bar Q)=2A-2\sum_i(R_i-\bar R)(Q_i-\bar Q).
\]
よって $\sum_i(R_i-\bar R)(Q_i-\bar Q)=A-D/2$ であり，
\[
r_S=\frac{A-D/2}{\sqrt{A\cdot A}}=1-\frac{D}{2A}=1-\frac{6D}{n(n^2-1)}.
\]

\ans{2}
以下，$(R_1,\ldots,R_n)$ を与えた条件の下で考える。$(Q_1,\ldots,Q_n)$ は $n!$ 通りの順列の上で一様に分布するので，
各 $Q_i$ は $\{1,\ldots,n\}$ 上の一様分布に従う。よって
\[
\E[Q_i]=\frac{n+1}2,\qquad
\Var(Q_i)=\frac{(n+1)(2n+1)}6-\frac{(n+1)^2}4=\boxed{\frac{n^2-1}{12}}.
\]
また $\sum_{i=1}^nQ_i=n(n+1)/2$ は定数であるから，$i\ne j$ に対する共分散を $\gamma$ とおくと（対称性より $i,j$ によらない）
\[
0=\Var\Bigl(\sum_{i=1}^nQ_i\Bigr)=n\cdot\frac{n^2-1}{12}+n(n-1)\gamma
\quad\Longrightarrow\quad
\Cov(Q_i,Q_j)=\gamma=\boxed{-\frac{n+1}{12}}.
\]
$a_i=R_i-\bar R$ とおくと $\sum_ia_i=0$，$\sum_ia_i^2=A$ であり，〔1〕の記号で
\[
r_S=\frac1A\sum_{i=1}^na_i(Q_i-\bar Q)=\frac1A\sum_{i=1}^na_iQ_i
\]
と書ける（$\sum_ia_i\bar Q=0$ を用いた）。$a_i$ は条件付けにより定数として扱えるので
\[
\E[r_S\mid R_1,\ldots,R_n]=\frac1A\cdot\frac{n+1}2\sum_ia_i=0 .
\]
分散は，$\sum_{i\ne j}a_ia_j=(\sum_ia_i)^2-\sum_ia_i^2=-A$ を用いて
\[
\Var\Bigl(\sum_ia_iQ_i\Bigm| R_1,\ldots,R_n\Bigr)=\frac{n^2-1}{12}\sum_ia_i^2+\gamma\sum_{i\ne j}a_ia_j
=\left(\frac{n^2-1}{12}+\frac{n+1}{12}\right)A=\frac{n(n+1)}{12}A .
\]
よって
\[
\Var(r_S\mid R_1,\ldots,R_n)=\frac{n(n+1)}{12A}=\frac{n(n+1)}{n(n^2-1)}=\frac1{n-1}.
\]
条件付き期待値・条件付き分散はいずれも $(R_1,\ldots,R_n)$ によらない定数であるから，
\[
\E[r_S]=\E\bigl[\E[r_S\mid R]\bigr]=0,\qquad
\Var(r_S)=\E\bigl[\Var(r_S\mid R)\bigr]+\Var\bigl(\E[r_S\mid R]\bigr)=\frac1{n-1}.
\]

\ans{3}
順位は次のとおりである。
\begin{center}
\begin{tabular}{cccccc}
\toprule
$i$ & 1 & 2 & 3 & 4 & 5\\
\midrule
$R_i$ & 3 & 1 & 5 & 2 & 4\\
$Q_i$ & 4 & 1 & 5 & 2 & 3\\
$R_i-Q_i$ & $-1$ & 0 & 0 & 0 & 1\\
\bottomrule
\end{tabular}
\end{center}
$D=2$ であるから，〔1〕より
\[
r_S=1-\frac{6\times2}{5\times24}=\boxed{0.9}.
\]
$r_S$ は $D$ の減少関数なので，$r_S\ge0.9$ は $D\le2$ と同値である。
$H_0$ の下で，$(R_i)$ を与えたとき $(Q_i)$ は $5!=120$ 通りの順列を等確率でとる。
$d_i=R_i-Q_i$ は整数で $\sum_id_i=0$ を満たす。
\begin{itemize}
\item $D=0$ となるのは，すべての $i$ で $Q_i=R_i$ となる1通り。
\item $D=2$ となるのは，$d_i$ のうちちょうど2つが $\pm1$（和が0なので $+1$ と $-1$ が1つずつ）で残りが0の場合である。
これは $R$ の順位が $k$ と $k+1$ である2個体の間で $Q$ の順位が入れ替わっている場合であり，$k=1,2,3,4$ の4通り。
\item $D=1$ となる場合はない（$d_i$ が1つだけ $\pm1$ では和が0にならない）。
\end{itemize}
よって正確な片側 $p$ 値は
\[
\Prob(r_S\ge0.9)=\Prob(D\le2)=\frac{1+4}{120}=\boxed{\frac1{24}\simeq0.042}.
\]
正規近似では $\sqrt{n-1}\,r_S=2\times0.9=1.80$ であり，付表1より $Q(1.80)=0.0359$ なので，片側 $p$ 値は約 $\boxed{0.036}$ である。
正規近似の $p$ 値は正確な値よりやや小さいが，どちらも0.05より小さく，有意水準5\%の片側検定では $H_0$ は棄却される。
$n=5$ では $r_S$ のとり得る値が少なく離散的であるため，近似の精度は高くない。

\ans{4}
$0<x<1$ に対して，$\int_0^1(1-2y)\,dy=0$ より
\[
h_X(x)=\int_0^1\{1+\theta(1-2x)(1-2y)\}\,dy=1 .
\]
同様に $h_Y(y)=1$（$0<y<1$）である。すなわち，$X_1$ と $Y_1$ の周辺分布はどちらも
\fbox{区間 $(0,1)$ 上の一様分布} である。
したがって $\E[X_1]=\E[Y_1]=1/2$，$\Var(X_1)=\Var(Y_1)=1/12$ である。また
\[
\E[X_1Y_1]=\int_0^1\!\!\int_0^1xy\,dx\,dy+\theta\left(\int_0^1x(1-2x)\,dx\right)^2
=\frac14+\theta\left(\frac12-\frac23\right)^2=\frac14+\frac\theta{36}.
\]
よって $\Cov(X_1,Y_1)=\theta/36$ であり，相関係数は
\[
\frac{\theta/36}{1/12}=\boxed{\frac\theta3}.
\]

\ans{5}
$n=2$ のとき，順位の組は $(R_1,R_2),(Q_1,Q_2)\in\{(1,2),(2,1)\}$ であり，〔1〕より $r_S=1-D$ で，
$X$ と $Y$ の大小の向きが一致すれば（$D=0$）$r_S=1$，逆であれば（$D=2$）$r_S=-1$ である。したがって
\[
\E[r_S]=\pi_c-(1-\pi_c)=2\pi_c-1,\qquad \pi_c=\Prob\bigl((X_1-X_2)(Y_1-Y_2)>0\bigr).
\]
2つの観測は独立同分布であるから，対称性より
\[
\pi_c=\Prob(X_1<X_2,\,Y_1<Y_2)+\Prob(X_1>X_2,\,Y_1>Y_2)=2\,\Prob(X_1<X_2,\,Y_1<Y_2).
\]
$(X_1,Y_1)$ の同時分布関数を $H(x,y)=\Prob(X_1\le x,Y_1\le y)$ とすると，$0<x<1,\ 0<y<1$ で
\[
H(x,y)=\int_0^x\!\!\int_0^y\{1+\theta(1-2s)(1-2t)\}\,dt\,ds=xy+\theta\,x(1-x)\,y(1-y)
\]
である（$\int_0^x(1-2s)\,ds=x(1-x)$）。$(X_2,Y_2)$ で条件付けると，独立性より
\[
\Prob(X_1<X_2,\,Y_1<Y_2)=\E[H(X_2,Y_2)]
=\int_0^1\!\!\int_0^1\{xy+\theta x(1-x)y(1-y)\}\{1+\theta(1-2x)(1-2y)\}\,dx\,dy .
\]
被積分関数を展開すると，各項は $x$ の関数と $y$ の関数の積であり，
\begin{gather*}
\int_0^1x\,dx=\frac12,\qquad
\int_0^1x(1-x)\,dx=\frac16,\qquad
\int_0^1x(1-2x)\,dx=-\frac16,\\
\int_0^1x(1-x)(1-2x)\,dx=\frac12-1+\frac12=0
\end{gather*}
を用いると
\[
\E[H(X_2,Y_2)]=\frac14+\theta\left(\frac16\right)^2+\theta\left(-\frac16\right)^2+\theta^2\cdot0^2=\frac14+\frac\theta{18}.
\]
よって $\pi_c=\dfrac12+\dfrac\theta9$ であり，
\[
\E[r_S]=2\pi_c-1=\boxed{\frac{2\theta}9}.
\]
（参考：この分布では周辺分布が一様分布なので，〔4〕の $\theta/3$ は母集団における順位相関に当たる。$n=2$ の $r_S$ の期待値 $2\theta/9$ はこれより絶対値が小さい。）
\newpage

\qhead{4}
\ans{1}
確率モデルは $X\sim\mathrm{Bin}(n,p)$ であり，確率関数は $\Prob(X=x)=\binom nxp^x(1-p)^{n-x}$（$x=0,1,\ldots,n$）である。
したがって尤度関数と対数尤度は
\[
L(p)=\binom nxp^x(1-p)^{n-x},\qquad
\ell(p)=\log\binom nx+x\log p+(n-x)\log(1-p).
\]
$M_1$ の下では $0<p<1$ で
\[
\ell'(p)=\frac xp-\frac{n-x}{1-p}=\frac{x-np}{p(1-p)}
\]
であり，$0<x<n$ より $\ell'(p)$ は $p<x/n$ で正，$p>x/n$ で負となる。よって $\ell(p)$ は $\hat p=x/n$ で最大となり，
$\hat\ell_1=\ell(x/n)$ である。$M_0$ には自由母数がないので $\hat\ell_0=\ell(1/2)=\log\binom nx-n\log2$ である。したがって
\begin{align*}
\Lambda=2(\hat\ell_1-\hat\ell_0)
&=\boxed{2\left\{x\log\frac xn+(n-x)\log\frac{n-x}n+n\log2\right\}}\\
&=2\left\{x\log\frac{2x}n+(n-x)\log\frac{2(n-x)}n\right\}.
\end{align*}
$\mathrm{AIC}=-2\hat\ell+2k$ であるから，$\mathrm{AIC}_0=-2\hat\ell_0$，$\mathrm{AIC}_1=-2\hat\ell_1+2$ であり，
\[
\mathrm{AIC}_1<\mathrm{AIC}_0\iff\boxed{\Lambda>2}.
\]
同様に $\mathrm{BIC}_0=-2\hat\ell_0$，$\mathrm{BIC}_1=-2\hat\ell_1+\log n$ であるから，
\[
\mathrm{BIC}_1<\mathrm{BIC}_0\iff\boxed{\Lambda>\log n}.
\]

\ans{2}
$2x/n=1.2$，$2(n-x)/n=0.8$ であるから $\Lambda=2(60\log1.2+40\log0.8)$ である。
付表5より $\log_{10}1.2=0.0792$，$\log_{10}0.8=-0.0969$ なので
\[
\log1.2=2.3026\times0.0792=0.1824,\qquad \log0.8=2.3026\times(-0.0969)=-0.2231 .
\]
\[
\Lambda=2\{60\times0.18237+40\times(-0.22312)\}=2(10.942-8.925)=2\times2.017=4.034,\qquad\boxed{\Lambda\simeq4.03}.
\]
（対数を小数第4位に丸めてから掛けると $4.04$ となる。この程度の差は許容する。正確な値は $4.027$。）
\begin{itemize}
\item AIC：$\Lambda>2$ なので \fbox{$M_1$ を選択} する。
\item BIC：$\log n=\log100=2\times2.3026=4.605$ であり，$\Lambda<\log n$ なので \fbox{$M_0$ を選択} する。
\item 尤度比検定：帰無仮説の下で $\Lambda$ は近似的に自由度 $1$（$M_1$ と $M_0$ の自由母数の個数の差）のカイ二乗分布に従う。
付表3より $\chi^2_{0.05}(1)=3.84$ であり，$\Lambda\simeq4.03>3.84$ なので，有意水準5\%で \fbox{帰無仮説 $p=1/2$ は棄却される}。
\end{itemize}

\ans{3}
$M_0$ の下で $\Lambda$ は自由度1のカイ二乗分布に分布収束する。この極限分布は連続であり，$Z\sim N(0,1)$ とすると $Z^2$ の分布に等しい。
AIC が $M_1$ を選択する事象は $\{\Lambda>2\}$ であるから，
\[
\lim_{n\to\infty}\Prob(\Lambda>2)=\Prob(Z^2>2)=2\,\Prob(Z>\sqrt2)=2Q(1.4142).
\]
付表1より $Q(1.41)=0.0793$，$Q(1.42)=0.0778$ であり，線形補間すると $Q(1.4142)\simeq0.0793-0.42\times0.0015=0.0787$ となる。よって
\[
\lim_{n\to\infty}\Prob(\text{AIC が }M_1\text{ を選択})\simeq2\times0.0787=\boxed{0.16}.
\]
すなわち，標本サイズをいくら大きくしても，AIC は約16\%の確率で誤って $M_1$ を選択する。

次に BIC が $M_1$ を選択する事象は $\{\Lambda>\log n\}$ である。任意の定数 $K>0$ をとると，$\log n>K$ となる十分大きな $n$ に対して
\[
\Prob(\Lambda>\log n)\le\Prob(\Lambda>K)
\]
が成り立つ。$n\to\infty$ とすると右辺は $\Prob(Z^2>K)$ に収束するから，
\[
\limsup_{n\to\infty}\Prob(\Lambda>\log n)\le\Prob(Z^2>K).
\]
$K$ は任意であり，$K\to\infty$ で $\Prob(Z^2>K)\to0$ であるから，$\Prob(\Lambda>\log n)\to0$ である。
すなわち BIC が誤って $M_1$ を選択する確率は0に収束する。

\ans{4}
$M_1$ の下では，$p$ を与えたときの $X$ の確率関数が $\binom nxp^x(1-p)^{n-x}$，$p$ の事前密度が $1$（$0<p<1$）である。
周辺確率関数は同時分布を $p$ について積分したものであり，ベータ関数の公式
$\int_0^1p^{a-1}(1-p)^{b-1}dp=\Gamma(a)\Gamma(b)/\Gamma(a+b)$ を用いると
\begin{align*}
m_1(x)&=\int_0^1\binom nxp^x(1-p)^{n-x}\,dp=\binom nx\frac{\Gamma(x+1)\Gamma(n-x+1)}{\Gamma(n+2)}\\
&=\frac{n!}{x!\,(n-x)!}\cdot\frac{x!\,(n-x)!}{(n+1)!}=\boxed{\frac{1}{n+1}}.
\end{align*}
（$x$ によらず，$0,1,\ldots,n$ の上の一様分布になる。）一方 $m_0(x)=\binom nx2^{-n}$ であるから
\[
B_{10}=\frac{m_1(x)}{m_0(x)}=\boxed{\frac{2^n}{(n+1)\binom nx}}.
\]

\ans{5}
$\hat\ell_1-\hat\ell_0=\Lambda/2$ より
\[
m_0(x)=\binom nx2^{-n}=\binom nx\hat p^{\,x}(1-\hat p)^{n-x}\,e^{-\Lambda/2}
\approx\frac{e^{-\Lambda/2}}{\sqrt{2\pi n\hat p(1-\hat p)}}
\]
である。したがって
\[
B_{10}=\frac{1}{(n+1)\,m_0(x)}\approx\frac{\sqrt{2\pi n\hat p(1-\hat p)}}{n+1}\,e^{\Lambda/2},
\]
\[
2\log B_{10}\approx\Lambda+\log n-2\log(n+1)+\log\{2\pi\hat p(1-\hat p)\}
\approx\boxed{\Lambda-\log n+\log\{2\pi\hat p(1-\hat p)\}}
\]
となる（$n$ が大きいとき $2\log(n+1)\approx2\log n$）。〔1〕より $\mathrm{BIC}_0-\mathrm{BIC}_1=\Lambda-\log n$ であるから，
\[
2\log B_{10}\approx(\mathrm{BIC}_0-\mathrm{BIC}_1)+\log\{2\pi\hat p(1-\hat p)\}
\]
であり，$2\log B_{10}$ と BIC の差は，$n$ が大きくなっても有界な項 $\log\{2\pi\hat p(1-\hat p)\}$ だけしか違わない。
すなわち BIC によるモデル選択は，ベイズファクターが1より大きいか小さいかによる選択を近似したものになっている。

$n=100$，$x=60$ のとき，$\hat p=0.6$，$2\pi n\hat p(1-\hat p)=2\times3.1416\times24=150.80$，$\sqrt{150.80}=12.28$ である。
〔2〕より $\Lambda/2=2.017$ であり，付表5の $e^{1.00}=2.7183$，$e^{0.01}=1.0101$，$e^{0.02}=1.0202$ を用いると
\[
e^{2.017}=e^{1.00}\times e^{1.00}\times e^{0.017}\simeq2.7183^2\times1.0173=7.3892\times1.0173=7.517
\]
（$e^{0.017}$ は線形補間による）。よって
\[
B_{10}\approx\frac{12.28\times7.517}{101}=\boxed{0.91}.
\]
（参考：$2^{100}/\{101\binom{100}{60}\}$ の正確な値は $0.913$ である。）

【考察】〔2〕では有意水準5\%の検定で $M_0$（$p=1/2$）が棄却され，AIC も $M_1$ を選んだ。
しかしベイズファクターは $B_{10}\simeq0.91<1$ であり，データは $M_1$ よりむしろ $M_0$ をわずかに支持している。これは BIC が $M_0$ を選んだことと整合する。
上の近似式から，$\Lambda$（すなわち検定の $p$ 値）を固定したまま $n$ を大きくすると，$2\log B_{10}\approx\Lambda-\log n+O(1)\to-\infty$ となり $B_{10}\to0$ である。
$M_1$ の事前分布は $p=1/2$ から離れた値にも広く確率を与えるので，$n$ が大きいときに $\hat p$ が $1/2$ のごく近くにあるデータは，
検定では有意であっても $M_1$ の下では平均的に起こりにくく，$M_0$ の方がよく説明する。
有意水準を固定した検定と，BIC・ベイズファクターとでは，$n$ が大きいときに結論が食い違うことがある。
\newpage

\qhead{5}
\ans{1}
$\E[X_1^2]=\E[X_3^2]=1$，$\E[X_1X_3]=\rho_{13}$ であるから
\[
\E[(X_1-aX_3)^2]=1-2a\rho_{13}+a^2=(a-\rho_{13})^2+1-\rho_{13}^2 .
\]
よって最小にする $a$ は $\boxed{a=\rho_{13}}$ であり，同様に $b=\rho_{23}$ である。すなわち
\[
e_1=X_1-\rho_{13}X_3,\qquad e_2=X_2-\rho_{23}X_3 .
\]
$(e_1,e_2,X_3)$ は3変量正規分布に従う $(X_1,X_2,X_3)$ の1次変換であるから，これも3変量正規分布に従う。また
\[
\Cov(e_1,X_3)=\Cov(X_1,X_3)-\rho_{13}\Var(X_3)=\rho_{13}-\rho_{13}=0,\qquad \Cov(e_2,X_3)=\rho_{23}-\rho_{23}=0 .
\]
多変量正規分布では，2つの部分ベクトルの間の共分散がすべて0であれば両者は独立である（同時密度が積に分解される）。
よって $(e_1,e_2)$ は $X_3$ と独立である。

\ans{2}
$\E[e_1]=\E[e_2]=0$ であり，
\begin{align*}
\Var(e_1)&=\Var(X_1)-2\rho_{13}\Cov(X_1,X_3)+\rho_{13}^2\Var(X_3)=1-\rho_{13}^2,\qquad \Var(e_2)=1-\rho_{23}^2,\\
\Cov(e_1,e_2)&=\Cov(X_1,X_2)-\rho_{23}\Cov(X_1,X_3)-\rho_{13}\Cov(X_3,X_2)+\rho_{13}\rho_{23}\Var(X_3)\\
&=\rho_{12}-\rho_{13}\rho_{23}-\rho_{13}\rho_{23}+\rho_{13}\rho_{23}=\rho_{12}-\rho_{13}\rho_{23}.
\end{align*}
したがって
\[
\rho_{12\cdot3}=\boxed{\frac{\rho_{12}-\rho_{13}\rho_{23}}{\sqrt{(1-\rho_{13}^2)(1-\rho_{23}^2)}}}.
\]

\ans{3}
$(X_1,X_2,X_3)$ の同時密度を $f(x_1,x_2,x_3)$，$X_3$ の周辺密度（標準正規密度）を $\varphi(x_3)$ とすると，
条件付き密度は定義により $f(x_1,x_2\mid x_3)=f(x_1,x_2,x_3)/\varphi(x_3)$ である。

〔1〕〔2〕より，$(e_1,e_2)$ は期待値 $(0,0)$，分散 $1-\rho_{13}^2,\ 1-\rho_{23}^2$，共分散 $\rho_{12}-\rho_{13}\rho_{23}$ の2変量正規分布に従う。
その密度を $g(u_1,u_2)$ とする。$(e_1,e_2)$ と $X_3$ は独立であるから，$(e_1,e_2,X_3)$ の同時密度は $g(u_1,u_2)\varphi(x_3)$ である。
変換 $(x_1,x_2,x_3)\mapsto(u_1,u_2,x_3)=(x_1-\rho_{13}x_3,\ x_2-\rho_{23}x_3,\ x_3)$ のヤコビアンは1であるから
\[
f(x_1,x_2,x_3)=g(x_1-\rho_{13}x_3,\ x_2-\rho_{23}x_3)\,\varphi(x_3),
\qquad
f(x_1,x_2\mid x_3)=g(x_1-\rho_{13}x_3,\ x_2-\rho_{23}x_3).
\]
これは $g$ を $(\rho_{13}x_3,\rho_{23}x_3)$ だけ平行移動した密度である。したがって，$X_3=x_3$ を与えたときの $(X_1,X_2)$ の条件付き分布は
\[
\boxed{\text{期待値 }(\rho_{13}x_3,\ \rho_{23}x_3),\ \text{分散 }1-\rho_{13}^2,\ 1-\rho_{23}^2,\ \text{共分散 }\rho_{12}-\rho_{13}\rho_{23}\ \text{の2変量正規分布}}
\]
である。分散・共分散は $x_3$ を含まないので，条件付き分布の相関係数は $x_3$ によらず
\[
\frac{\rho_{12}-\rho_{13}\rho_{23}}{\sqrt{(1-\rho_{13}^2)(1-\rho_{23}^2)}}=\rho_{12\cdot3}
\]
である。

\ans{4}
(a) $\rho_{12}-\rho_{13}\rho_{23}=0.48-0.6\times0.8=0$ より $\boxed{\rho_{12\cdot3}=0}$。
$X_3$ を考慮しなければ $X_1$ と $X_2$ には正の相関（0.48）があるが，$X_3$ の値を固定すると無相関であり，
〔3〕より条件付き分布は2変量正規分布なので，$X_3$ を与えた下で $X_1$ と $X_2$ は独立である。
すなわち，$X_1$ と $X_2$ の相関はすべて共通の変数 $X_3$ との関連によって生じた見かけ上のものである。

(b) $\rho_{12\cdot3}=\dfrac{0-0.6\times0.6}{\sqrt{(1-0.36)(1-0.36)}}=-\dfrac{0.36}{0.64}=\boxed{-0.5625}$。
$X_3$ を考慮しなければ $X_1$ と $X_2$ は無相関であり，2変量正規分布に従うので独立である。
しかし $X_3$ の値を固定すると負の相関（$-0.5625$）が現れる。周辺では独立な2変数でも，両方と関連する変数で条件付けると従属になる。

（いずれの場合も行列式 $1+2\rho_{12}\rho_{13}\rho_{23}-\rho_{12}^2-\rho_{13}^2-\rho_{23}^2$ は (a) $0.2304$，(b) $0.28$ で正であり，分散共分散行列は正則である。）

\ans{5}
$X_1=\rho_{13}X_3+e_1$，$X_2=\rho_{23}X_3+e_2$ を代入すると
\[
X_1-\beta_2X_2-\beta_3X_3=(e_1-\beta_2e_2)+(\rho_{13}-\beta_2\rho_{23}-\beta_3)X_3 .
\]
$(e_1,e_2)$ と $X_3$ は独立で期待値が0，$\E[X_3^2]=1$ であるから，交差項の期待値は0となり
\[
\E[(X_1-\beta_2X_2-\beta_3X_3)^2]=\E[(e_1-\beta_2e_2)^2]+(\rho_{13}-\beta_2\rho_{23}-\beta_3)^2 .
\]
第2項は，どの $\beta_2$ に対しても $\beta_3=\rho_{13}-\beta_2\rho_{23}$ ととれば0にできる。第1項は
\[
\E[(e_1-\beta_2e_2)^2]=\Var(e_2)\left\{\beta_2-\frac{\Cov(e_1,e_2)}{\Var(e_2)}\right\}^2+\Var(e_1)-\frac{\Cov(e_1,e_2)^2}{\Var(e_2)}
\]
と平方完成できるので，最小にする $\beta_2$ は
\[
\beta_2=\frac{\Cov(e_1,e_2)}{\Var(e_2)}=\boxed{\frac{\rho_{12}-\rho_{13}\rho_{23}}{1-\rho_{23}^2}}
=\rho_{12\cdot3}\sqrt{\frac{1-\rho_{13}^2}{1-\rho_{23}^2}}
\]
である。$\sqrt{(1-\rho_{13}^2)/(1-\rho_{23}^2)}>0$ であるから，$\beta_2$ の符号は $\rho_{12\cdot3}$ の符号と一致する（一方が0なら他方も0）。

最小値は，$\Cov(e_1,e_2)^2=\rho_{12\cdot3}^2\Var(e_1)\Var(e_2)$ を用いて
\[
\Var(e_1)-\frac{\Cov(e_1,e_2)^2}{\Var(e_2)}=\Var(e_1)(1-\rho_{12\cdot3}^2)=(1-\rho_{13}^2)(1-\rho_{12\cdot3}^2)
\]
となる。

（参考：$X_1$ を $X_2,X_3$ で予測するときの重相関係数を $R$ とすると，この最小値は $1-R^2$ であるから，
$1-R^2=(1-\rho_{13}^2)(1-\rho_{12\cdot3}^2)$ が成り立つ。$X_3$ だけによる予測の残差分散 $1-\rho_{13}^2$ が，$X_2$ を加えることで $1-\rho_{12\cdot3}^2$ 倍になる。）

\ans{6}
$\rho_{12}$ がとり得る値は，分散共分散行列（各成分の分散が1なので相関行列）$\Sigma$ が正則な分散共分散行列，すなわち正定値となるような値である。
$(e_1,e_2,X_3)$ は
\[
\begin{pmatrix}e_1\\e_2\\X_3\end{pmatrix}=A\begin{pmatrix}X_1\\X_2\\X_3\end{pmatrix},\qquad
A=\begin{pmatrix}1&0&-\rho_{13}\\0&1&-\rho_{23}\\0&0&1\end{pmatrix}
\]
と書け，$A$ は正則（$\det A=1$）であるから，$\Sigma$ が正定値であることと，$(e_1,e_2,X_3)$ の分散共分散行列 $A\Sigma A^\top$ が正定値であることは同値である。
〔1〕〔2〕の分散・共分散の計算は $\Sigma$ の正則性を使っていないので，
\[
A\Sigma A^\top=\begin{pmatrix}1-\rho_{13}^2&\rho_{12}-\rho_{13}\rho_{23}&0\\ \rho_{12}-\rho_{13}\rho_{23}&1-\rho_{23}^2&0\\0&0&1\end{pmatrix}
\]
である。これが正定値となるのは，左上の $2\times2$ 行列が正定値のとき，すなわち対角成分が正で行列式が正のときである。
$|\rho_{13}|<1$，$|\rho_{23}|<1$ より対角成分は正であるから，条件は
\[
(\rho_{12}-\rho_{13}\rho_{23})^2<(1-\rho_{13}^2)(1-\rho_{23}^2)
\]
となる（これは $|\rho_{12\cdot3}|<1$ と同じ条件である）。$\rho_{13}=0.6$，$\rho_{23}=0.8$ を代入すると
\[
(\rho_{12}-0.48)^2<0.64\times0.36=0.2304=0.48^2
\quad\Longleftrightarrow\quad\boxed{0<\rho_{12}<0.96}.
\]
（参考：$\det\Sigma=1+2\rho_{12}\rho_{13}\rho_{23}-\rho_{12}^2-\rho_{13}^2-\rho_{23}^2=(1-\rho_{13}^2)(1-\rho_{23}^2)(1-\rho_{12\cdot3}^2)$ が成り立つ。
$X_1$ と $X_2$ がともに $X_3$ と強い正の相関をもつとき，$X_1$ と $X_2$ の相関は0以下にはなれない。）
\end{document}
